% 张三慧大学物理学（第4版）V2第26章《波粒二象性》
% 主教材扫描PDF物理页368--369（印刷351--352）；配套解答S PDF354--360。
% 22道原编号章末习题；思考题另编，不把26.1--26.13思考题冒充章末习题。
\section*{第26章\quad 波粒二象性}
\subsection*{章末习题26.1--26.22}
\paragraph{习题 26.1}
晴天夜晚地面降温主要由于地面热辐射。
若地面温度 \(-5^\circ\mathrm C\)，按黑体辐射
估算每 \(1\,\mathrm{m^2}\) 地面向外发射的总功率。
\paragraph{习题 26.2}
太阳黑体辐射的光谱辐射出射度
\(M_\nu\) 极大处频率
\(\nu_m=3.4\times10^{14}\,\mathrm{Hz}\)。
求(1)太阳表面温度 \(T\)；(2)太阳表面黑体总出射度
\(M_0\)。须对\emph{频率谱}的峰值使用相应Wien常数。
\paragraph{习题 26.3}
地球表面太阳辐照强度
\(1.0\times10^3\,\mathrm{W/m^2}\)。
太阳能水箱黑色受光面正对阳光。
用单面黑体热平衡模型估算水箱的最终水温
（摄氏度），忽略其他各面热辐射。
\paragraph{习题 26.4}
太阳总辐射功率
\(P_\odot=3.9\times10^{26}\,\mathrm W\)。
(1) 对绕太阳作半径 \(r\) 的圆轨道的行星，
在黑体吸收/辐射热平衡下证明
\[
T^4=\frac{P_\odot}{16\pi\sigma r^2}.
\]
(2) 地球轨道半径
\(r_E=1.5\times10^{11}\,\mathrm m\)，
冥王星轨道半径
\(r_P=5.9\times10^{12}\,\mathrm m\)，
求各自有效黑体表面温度。
\paragraph{习题 26.5}
Procyon B 距地球11光年，在地面接收到的
辐射强度 \(I=1.7\times10^{-12}\,\mathrm{W/m^2}\)，
该星表面温度6600K。按黑体发射与各向同性传播
估算恒星半径。
\paragraph{习题 26.6}
宇宙微波背景辐射相当于温度3K的
各向同性黑体辐射。
(1) 其\emph{频率谱}辐射出射度 \(M_\nu\)
在哪一频率处最大？(2) 估算地球表面从各方向
接收这种辐射的总功率，取地球半径的标准近似值。
\paragraph{习题 26.7}
由普朗克黑体\emph{频率谱}公式
\[
M_\nu=\frac{2\pi h\nu^3}{c^2}
\frac{1}{\mathrm e^{h\nu/(k_BT)}-1}
\]
导出按\emph{波长}表示的谱出射度
\[
M_\lambda=\frac{2\pi hc^2}{\lambda^5}
\frac{1}{\mathrm e^{hc/(\lambda k_BT)}-1}.
\]
\paragraph{习题 26.8}
令 \(w_\nu\) 为黑体空腔内辐射的频谱能量密度；
证明单位面积黑体辐射频谱出射度
\(\boxed{M_\nu=(c/4)w_\nu}\)。
\paragraph{习题 26.9}
对普朗克式 \(M_\nu\) 求导，证明频率谱峰值
满足 \(\nu_m=C_\nu T\)（Wien位移定律），
可不求无量纲方程的根数值。
\paragraph{习题 26.10}
通过对普朗克辐射出射度谱积分，
证明黑体总出射度 \(M=\sigma T^4\)
（Stefan--Boltzmann定律）；可不求常数 \(\sigma\) 数值。
\paragraph{习题 26.11}
铝的逸出功 \(A=4.2\,\mathrm{eV}\)。
用波长 \(\lambda=200\,\mathrm{nm}\) 的光照射，
求(1)光电子最大动能；(2)遏止电压；
(3)铝的红限波长。
\paragraph{习题 26.12}
银河系恒星光的平均辐射能量密度
\(u=10^{-15}\,\mathrm{J/m^3}\)，
光子平均波长假设为500nm。
求光子数密度。
\paragraph{习题 26.13}
功率1.0W的灯泡在距其1.0m的
钾片上产生光电效应（钾逸出功2.25eV）。
若一个电子只能接收到半径
\(r=1.3\times10^{-10}\,\mathrm m\) 的
圆面积上的辐射能量，按经典连续能量积累模型
估算这个电子积蓄逸出功所需时间。
指出与真实光电效应的时间尺度的差别。
\paragraph{习题 26.14}
在实验室系中一光子能量5eV，
一质子沿光子的传播方向以 \(c/2\) 运动。
求质子静止参考系中光子的能量。
\paragraph{习题 26.15}
能量 \(0.60\,\mathrm{MeV}\) 的入射X射线光子
被自由电子散射后波长增长20\%。
求反冲电子获得的动能。
\paragraph{习题 26.16}
一个能量 \(4.0\times10^3\,\mathrm{eV}\)
的光子与初始静止的自由电子碰撞。
求电子所能得到的最大动能。
\paragraph{习题 26.17}
仅用动量守恒和能量守恒证明：
单独的自由电子不能一次完全吸收一个光子。
\paragraph{习题 26.18}
能量 \(5.0\times10^4\,\mathrm{eV}\) 的光子
与一动能为 \(2.0\times10^4\,\mathrm{eV}\)
的电子作正碰，碰后光子向后折回。
求碰后光子与电子各自能量，
注意电子静止能量与动能区别。
\paragraph{习题 26.19}
电子和光子德布罗意波长均为
\(\lambda=0.20\,\mathrm{nm}\)。
求二者的动量和能量（电子须分别给总能量及动能）。
\paragraph{习题 26.20}
室温300K下的中子称为热中子。
按热运动的平均平动能量模型估算
中子的德布罗意波长，并注明
使用 \(E=(3/2)k_BT\) 的约定。
\paragraph{习题 26.21}
电子显微镜将电子经加速电压40kV加速。
求电子的德布罗意波长，说明相对论修正。
\paragraph{习题 26.22}
从德布罗意关系
\(\lambda=h/p\)、\(\nu=E/h\)
及相对论质量关系 \(E^2=p^2c^2+m_0^2c^4\)
出发，证明自由粒子的德布罗意波群群速
\(v_g=d\omega/dk\) 恰等于该粒子运动速度 \(v\)。
